\section*{章末問題}
\begin{enumerate}[label=\textbf{\thechapter.\arabic*},leftmargin=*]
\item\label{問:逆三角関数の枝}
$\arcsin(\sin\frac{4\pi}3)$，$\arccos(\cos\frac{5\pi}4)$，$\arctan(\tan(-\frac{2\pi}3))$を求めよ。また$y=\arcsin(x^2)$を$-1<x<1$で微分せよ。
\item\label{問:逆三角関数の積分練習}
$\int\frac1{x^2+2x+5}\dd x$と$\int_0^1\frac1{\sqrt{4-x^2}}\dd x$を求めよ。
\item\label{問:余接と正接の和}
定義から$(\cot x)'=-\csc^2x$を示せ。さらに$\frac1{\sin x\cos x}=\tan x+\cot x$を使って，$\int\frac1{\sin x\cos x}\dd x$を求め，答えを用いる区間の条件を述べよ。
\item\label{問:半角置換の区間分割}
半角置換を使って$\int_0^{2\pi}\frac1{2+\sin x}\dd x$を求めよ。$x=\pi$の前後を分け，置換後の無限の端点は極限として扱うこと。
\item\label{問:円の点とピタゴラス数}
円上の点$\left(\frac{12}{13},\frac5{13}\right)$に対応する傾き$t$を求めよ。$t=\frac nm$から式\eqref{式:ピタゴラス数}で作られる三整数を求め，共通因数で割れ。また$m,n$がともに奇数なら，三整数がすべて偶数になることを示せ。
\item\label{問:双曲線二倍角と積分}
指数関数の定義から$\sinh(2u)=2\sinh u\cosh u$を示せ。さらに$\int_0^{\log2}\cosh^2u\dd u$を求めよ。
\item\label{問:双曲線余弦の置換}
$x>1$で$\int\frac1{\sqrt{x^2-1}}\dd x$を，$x=\cosh u$によって求めよ。$u$の範囲を指定し，元の変数へ戻す対数の式も導け。
\item\label{問:双曲線置換と代数置換}
$\int_0^2\sqrt{4+x^2}\dd x$を，$x=2\sinh u$と$t=x+\sqrt{4+x^2}$の二通りの置換で求めよ。また二つの変数の間に$t=2\e^u$が成り立つことを確かめよ。
\end{enumerate}

\smallskip\noindent 解答は\pageref{解答:第10章}ページ。
